\documentclass[12pt,a4paper]{article}
\usepackage[utf8]{inputenc}
\usepackage[french]{babel}
\usepackage[T1]{fontenc}
\usepackage{amsmath}
\usepackage{amsfonts}
\usepackage{amssymb}
\usepackage[pdftex]{graphicx}
\usepackage{comment}
\usepackage{siunitx}
\usepackage[colorinlistoftodos,prependcaption,textsize=tiny]{todonotes}
\usepackage[left=2cm,right=2cm,top=2cm,bottom=2cm]{geometry}
\author{Angelo Ortiz}
\title{Rapport du Projet}

\begin{document}
\begin{titlepage}
  \centering
  \includegraphics[width=0.30\textwidth]{logo.jpg}\par\vspace{1cm}
  {\scshape\LARGE Sorbonne Universit\'e\par}
  \vspace{1cm}
  {\scshape\Large 2I013 : Projet (application)\par}
  \vspace{1.5cm}
  {\Large \bfseries Sujet :\par}
  {\huge\bfseries IA Football\par}
  \vspace{2cm}
  {\Large\itshape Fangzhou Ye\par}
  {\Large\itshape Angelo Ortiz\par}
  \vfill
  
  % Bottom of the page
  {\large Licence d'Informatique\par}
  {\large Ann\'ee 2017/2018\par}
\end{titlepage}
 
%\newpage
\tableofcontents
  
\newpage
  
\part*{Introduction}
\section*{Pr\'esentation}
Le projet a consist\'e en le d\'eveloppement d'intelligences artificielles de 
joueurs de football. Le simulateur du jeu \'etant d\'ej\`a fourni, notre 
travail a \'et\'e d'impl\'ementer des strat\'egies de jeu pour mieux r\'eussir 
les matches.
L’objectif de ce travail a \'et\'e d'apprendre \`a bien mener un projet 
collaboratif long sur un nouvel environnement de d\'eveloppement, \`a savoir 
Python.
\section*{Aper\c{c}u}
Voici un aper\c{c}u du jeu...
...
\section*{Arborescence}
Le projet a \'et\'e organis\'e en un module {\bfseries ia} dirigeant le 
comportement de nos joueurs et en divers fichiers de test.
\subsubsection*{ia/tools}
Ce fichier contient une classe repr\'esentant l'\'etat du jeu qui comporte des 
m\'ethodes nous facilitant l'acc\`es aux diverses carat\'eristiques de cet 
\'etat, ainsi que plusieurs fonctions utilitaires.

\subsubsection*{ia/conditions}
Comme le nom l'indique, ce fichier a \'et\'e d\'edi\'e aux diverses conditions 
qui d\'eclenchent les diff\'erents types de comportement.

\subsubsection*{ia/actions}
Ce fichier comporte les actions de base que l'on peut r\'ealiser \`a chaque 
instant du jeu.

\subsubsection*{ia/behaviour}
En ce qui concerne ce fichier, ici se trouvent les actions complexes, i.e.\ 
compos\'ees de plusieurs actions de base. En effet, pour une certaine phase de 
jeu, un joueur peut prendre diff\'erentes d\'ecisions.

\subsubsection*{ia/strategies}
C'est dans ce fichier que les strat\'egies ont \'et\'e r\'eunies. Ces 
strat\'egies dirigent le comportement des joueurs tout au long d'un 
match.

\subsubsection*{ia/gene\_optimisation}
Ce fichier contient les classes n\'ecessaires pour la mise en \oe uvre d'une 
impl\'ementation des algorithmes g\'en\'etiques.

\subsubsection*{ia/gs\_optimisation}
On retrouve ici le code utilis\'e pour suivre la m\'ethode de recherche en 
grille.

\subsubsection*{ia/}
init ? ML\_strats ? DT\_strats ? 
\todo{Il reste encore des fichiers}

\addcontentsline{toc}{part}{Introduction}
\newpage

\part{Architecture logicielle} %et demarche
Dans cette premi\`ere partie, on vous expliquera les choix faits 
concernant l'architecture logicielle.

\newpage

\part{Strat\'egies}
Dans cette deuxi\`eme partie, on vous d\'etaille le comportement de nos 
diff\'erents joueurs.

Compte tenu de la taille du terrain de jeu, le nombre maximum de joueurs par 
\'equipe est limit\'e \`a quatre. Ainsi, il y a trois cat\'egories de matches 
selon le nombre de joueurs par \'equipe : un, deux ou quatre.

\section*{Un joueur}
Pour l'\'equipe compos\'ee uniquement d'un joueur, on a d\'evelopp\'e un seul 
joueur que l'on appelle {\itshape Fonceur}. Au tout d\'ebut du 
projet, il faisait tout au maximum : il doit 
s'approcher de la balle \`a toute vitesse lorsqu'il ne la contr\^ole pas et il 
frappe avec toute sa puissance dans l'autre cas. On a d\'ecid\'e de garder le 
m\^eme nom, m\^eme si son comportement diff\`ere. En effet, \`a pr\'esent il 
avance avec la balle de petites distances et lorsqu'un adversaire lui fait 
opposition il essaie de le dribbler ; il frappe vers la cage adverse s'il se 
trouve dans la surface de r\'eparation. 

\section*{Deux joueurs}
Cette cat\'egorie a \'et\'e celle sur laquelle on a consacr\'e la plupart du 
projet. On compte trois joueurs fonctionnels que l'on appelle {\itshape 
  Attaquant}, {\itshape Gardien} et {\itshape D\'efenseur}. Les deux premiers 
pr\'esentent le m\^eme comportement lorsqu'ils se trouvent en situation 
d'attaque : \todo{reprendre ICI}

\newpage

\part{M\'ethodes d'optimisation}
Dans cette troisi\`eme partie, on vous d\'etaille les diff\'erents algorithmes 
appris tout au long du semestre pour l'am\'elioration des strat\'egies 
initialement propos\'ees.

\section{Recherche en grille}
...

\section{Algorithmes g\'en\'etiques}

\section{Arbres de d\'ecision}

\section{Apprentissage automatique}
Dans la derni\`ere partie du semestre, on s'est concentr\'e sur l'apprentissage 
automatique. On nous a pr\'esent\'e les deux grandes classes 
d'algorithmes d'apprentissage : supervis\'es et non supervis\'es, ainsi que 
l'apprentissage par renforcement. En effet, ce dernier type d'apprentissage 
impl\'ement\'e par l'algorithme de \texttt{Q-learning} a \'et\'e mise en place.

\newpage

\part*{Conclusion}
...
\addcontentsline{toc}{part}{Conclusion}

\end{document}
